\section{Conclusiones y Trabajo Extra}

\subsection{Conclusiones}
\begin{enumerate}[leftmargin=*]
    \item \textbf{Gestión Eficiente del Flujo Continuo:} El diseño del tablero Kanban en Jira (\textit{Tablero ID: 137}) eliminó las restricciones del modelo Scrum previo, permitiendo gestionar el inventario de 42 incidencias del proyecto SIGC-CUSCO con transparencia total.
    \item \textbf{Control de Capacidad mediante WIP:} La restricción de $\text{WIP} = 3$ en \textit{En curso} y $\text{WIP} = 2$ en \textit{En revisión} garantizó que los desarrolladores se concentren en cerrar tareas antes de abrir nuevas, reduciendo el \textit{Cycle Time} promedio a 2.1 días.
    \item \textbf{Resolución Rápida de Cuellos de Botella:} La técnica de \textit{Swarming} demostró ser un protocolo eficaz para descongestionar fases de revisión de código cuando se supera la capacidad máxima de la columna.
\end{enumerate}

\subsection{Trabajo Extra: Mantener el Tablero Actualizado durante la Semana}
En cumplimiento de la \textbf{Actividad 6 (Trabajo extra)} de la guía:
\begin{itemize}[leftmargin=*]
    \item \textbf{Sincronización Diaria (10 min):} Reunión matutina frente al tablero Jira para verificar que ninguna columna supere su WIP.
    \item \textbf{Política de Jalado (\textit{Pull}):} Ningún integrante jala tareas nuevas si la fase de revisión está al límite.
    \item \textbf{Revisión de Métricas Semanal:} Los viernes se audita el CFD en Jira para detectar variaciones en el tiempo de entrega y afinar la capacidad del equipo con la MDSJ.
\end{itemize}
